\documentclass[header={Ancillary Note: Weak vs Weak-*},notheorems]{study-note}

\begin{document}

\StudyTitleBlock{N6}{Ancillary Note}{Weak convergence vs. weak-* convergence}

\vspace{0.8em}

\begin{focusbox}
Weak and weak-* convergence are easy to confuse because the definitions look almost identical. The real difference is: which functionals are allowed to test convergence.
\end{focusbox}

\tableofcontents

\section{Weak Convergence}

\begin{defbox}{Weak convergence in $X$}
Let $X$ be a normed space. A sequence $(x_n)$ converges \textbf{weakly} to $x\in X$, written
\[
x_n\rightharpoonup x,
\]
if
\[
\phi(x_n)\to \phi(x)
\qquad \forall \phi\in X^*.
\]
\end{defbox}

\section{Weak-* Convergence}

\begin{defbox}{Weak-* convergence in $X^*$}
A sequence $(\phi_n)$ in $X^*$ converges \textbf{weak-*} to $\phi\in X^*$, written
\[
\phi_n \overset{*}{\rightharpoonup} \phi,
\]
if
\[
\phi_n(x)\to \phi(x)
\qquad \forall x\in X.
\]
\end{defbox}

\begin{warningbox}{Main difference}
Weak convergence in $X^*$ tests against all elements of $X^{**}$.

Weak-* convergence in $X^*$ tests only against elements of $X$.

So weak-* convergence is generally a weaker notion.
\end{warningbox}

\section{Why The Definitions Differ}

\begin{focusbox}
For elements of $X$, weak convergence is tested by functionals from $X^*$.

For elements of $X^*$, weak convergence is tested by functionals from $X^{**}$, while weak-* convergence is tested only by the canonical copy of $X$ inside $X^{**}$.
\end{focusbox}

\section{Relations Between Them}

\begin{thmbox}{Implication}
If $\phi_n\rightharpoonup \phi$ weakly in $X^*$, then
\[
\phi_n \overset{*}{\rightharpoonup} \phi.
\]
\end{thmbox}

\paragraph{Reason.}
Every $x\in X$ defines an element of $X^{**}$ by the canonical embedding, so weak convergence in $X^*$ gives the weak-* testing identities automatically.

\begin{warningbox}{Converse fails in general}
Weak-* convergence does not in general imply weak convergence in $X^*$.
\end{warningbox}

\section{Why Weak-* Is Useful}

\begin{thmbox}{Banach--Alaoglu idea}
The closed unit ball of $X^*$ is compact in the weak-* topology.
\end{thmbox}

\paragraph{Meaning.}
Weak-* convergence is designed so that bounded sets in the dual space have compactness properties that are often false for the weak topology.

\section{Hilbert-Space Simplification}

\begin{exbox}{Hilbert spaces}
If $H$ is a Hilbert space, then $H\cong H^*$ by Riesz representation. This makes weak convergence easier to visualize:
\[
x_n\rightharpoonup x
\Longleftrightarrow
(x_n,y)_H\to (x,y)_H
\quad \forall y\in H.
\]
\end{exbox}

\section{What To Remember Quickly}

\begin{focusbox}
Quick summary:
\begin{enumerate}
\item weak convergence in $X$ is tested by $X^*$;
\item weak convergence in $X^*$ is tested by $X^{**}$;
\item weak-* convergence in $X^*$ is tested only by $X$;
\item weak convergence implies weak-* convergence in the dual;
\item Banach--Alaoglu is the main reason weak-* convergence matters.
\end{enumerate}
\end{focusbox}

\end{document}
